Take an orthogonal standard basis, $a = e_i$ and $X = e_K$ with $|K| = k$ (both sides are bilinear). By \claimref{std-algebra.orthogonal-anticommute}, $e_K I = c\, e_{K^{c}}$ with $c$ a sign times $\prod_{j \in K} Q(e_j)$ and $K^{c}$ the complement, so $X I \in \Cl^{n-k}$ (\claimref{std-algebra.grade-decomposition}). If $i \in K$: $e_i \wedge e_K = 0$, and $e_i \notin K^{c}$ gives $e_i \rfloor e_{K^{c}} = 0$, so the first identity reads $0 = 0$; also $e_i \rfloor e_K = e_i e_K$ (\claimref{std-algebra.product-decomposition}) and $e_i \wedge (e_K I) = e_i e_K I$, since $e_i e_{K^c}$ has the single grade $n - k + 1$. If $i \notin K$: $e_i \wedge e_K = e_i e_K$, and $e_i e_K I = c\, e_i e_{K^{c}}$ with $i \in K^{c}$, which has the single grade $n - k - 1$, so equals $e_i \rfloor (e_K I)$; and both sides of the second identity vanish. This holds for every signature.

If $r = 0$, $I \tilde{I} = (-1)^{q} \ne 0$ (\claimref{std-algebra.pseudoscalar-commutation}), so $I$ is invertible and $X \mapsto XI$ is injective, hence an isomorphism $\Cl^{k} \to \Cl^{n-k}$ by $\binom{n}{k} = \binom{n}{n-k}$. For a nonzero blade $A$: $x \wedge (A I) = (x \rfloor A) I$, which vanishes exactly when $x \rfloor A = 0$. By \claimref{std-algebra.contraction-expansion} with $A = v_1 \wedge \dots \wedge v_k$, $x \rfloor A = \sum_i (-1)^{i-1} B(x, v_i)\, v_1 \wedge \dots \check{v}_i \dots \wedge v_k$, whose blades are independent when the $v_i$ are (\claimref{std-algebra.blade-subspace}); so $x \rfloor A = 0$ exactly when $x \perp [A]$. Thus $\{x : x \wedge AI = 0\} = [A]^{\perp}$, of dimension $n - k$ since $Q$ is nondegenerate; $AI \ne 0$ is an $(n-k)$-vector, so it is a blade with $[AI] = [A]^{\perp}$ by part (iii) of \claimref{std-algebra.blade-subspace}.
